% This is text associated with the open textbook "Open Electromagnetics"
% See immediately below for author, date
% This work is licensed under the Creative Commons Attribution-ShareAlike 4.0 International License. (CC BY SA 4.0) 
% To view a copy of the license, visit http://creativecommons.org/licenses/by-sa/4.0/

%ORB TITLE Attenuation Rate
%ORB AUTHOR 0        # S.W. Ellingson, Virginia Tech (ellingson@vt.edu)
%ORB DATE 20191115   # yyyymmdd
%ORB ABSTRACT

\label{m0154_Attenuation_Rate} % <-- Must match name of this file and module directory name.  Don't mess with this!
\index{attenuation rate|(}

\emph{Attenuation rate} is a convenient way to quantify loss in general media, including transmission lines, using the decibel scale.

Consider a transmission line carrying a wave in the $+z$ direction.
Let $P_0$ be the power at $z=0$.
Let $P_1$ be the power at $z=l$.
Then the power at $z=0$ relative to the power at $z=l$ is:
%
\begin{equation}
\frac{P_0}{P_1} = \frac{e^{-2\alpha\cdot 0}}{e^{-2\alpha\cdot l}} = e^{2\alpha l} ~~~\mbox{(linear units)}
\end{equation}
%
where $\alpha$ is the attenuation constant; that is, the real part of the propagation constant $\gamma=\alpha+j\beta$.
\index{propagation constant!attenuation}
Expressed in this manner, the power ratio is a loss; that is, a number greater than $1$ represents attenuation.
In the decibel scale, the loss is
%
\begin{flalign}
10\log_{10}\frac{P_0}{P_1} &= 10\log_{10} e^{2\alpha l} \nonumber \\
                           &= 20\alpha l \log_{10} e \nonumber \\
                           &\cong 8.69\alpha l~~\mbox{dB} 
\end{flalign}
%
Attenuation rate is defined as this quantity per unit length.  Dividing by $l$, we obtain:
%
\begin{equation}
\mbox{attenuation rate} \cong 8.69\alpha
\end{equation}
% 
This has units of dB/length, where the units of length are the same length units in which $\alpha$ is expressed.
For example, if $\alpha$ is expressed in units of m$^{-1}$, then attenuation rate has units of dB/m.
%
%%%%%%%%%%%%%%%%%%%%%%%%%%%%%%%%%%%%%%%%%%%%%%%
%ORB HIGHLIGHT
\begin{mdframed}[backgroundcolor=yellow!20,nobreak=true] 
Attenuation rate $\cong 8.69\alpha$ is the loss in dB, per unit length.
\end{mdframed}
%%%%%%%%%%%%%%%%%%%%%%%%%%%%%%%%%%%%%%%%%%%%%%%
%
The utility of the attenuation rate concept is that it allows us to quickly calculate loss for any distance of wave travel:
This loss is simply attenuation rate (dB/m) times length (m), which yields loss in dB.

%%%%%%%%%%%%%%%%%%%%%%%%%%%%%%%%%%%%%%%%%%%%%%%
%ORB EXAMPLE
~\\ \begin{example} 
\begin{mdframed}[backgroundcolor=brown!20,linecolor=brown!20]\raggedright
Attenuation rate in a long cable.
\\~\\
A particular coaxial cable has an attenuation constant $\alpha\cong 8.5\times 10^{-3}$~m$^{-1}$.
What is the attenuation rate and the loss in dB for 100~m of this cable?

The attenuation rate is
$$ \cong 8.69\alpha \cong \underline{0.0738~\mbox{dB/m}} $$
The loss in 100~m of this cable is
$$ \cong \left(0.0738~\mbox{dB/m}\right)\left(100~\mbox{m}\right) \cong \underline{7.4~\mbox{dB}} $$
Note that it would be entirely appropriate, and equivalent, to state that the attenuation rate for this cable is $7.4$~dB/(100~m).

\end{mdframed}
% note figures will need to go between "\end{mdframed}" and "\end{example}"
\end{example}
%%%%%%%%%%%%%%%%%%%%%%%%%%%%%%%%%%%%%%%%%%%%%%%

The concept of attenuation rate is used in precisely the same way to relate ratios of spatial power densities for unguided waves.  
This works because spatial power density has SI base units of W/m$^2$, so the common units of m$^{-2}$ in the numerator and denominator cancel in the power density ratio, leaving a simple power ratio.

%%%%%%%%%%%%%%%%%%%%%

%{\bf Additional Reading:}
%\begin{itemize}
%\item ``\href{https://en.wikipedia.org/wiki/Attenuation}{Attenuation}'' on Wikipedia.
%\end{itemize}

\index{attenuation rate|)}


